\chapter{The USSR and the West (1923--1927)}

The progress towards the establishment of normal relations with the western Powers, which had followed the introduction of NEP, suffered a setback during the turbulent year 1923. The year opened with the French occupation of the Ruhr in reprisal for a German default on reparations. In Great Britain, the fall of Lloyd George left Curzon in undivided control of foreign policy. In France, the no less inflexible Poincare was at the height of his power. In May 1923 a number of British protests against Soviet misdemeanours culminated in what came to be known as the ``Curzon ultimatum''. This rehearsed at length the activities of Soviet agents in Persia, Afghanistan and India in violation of the undertaking given in the Anglo-Soviet trade agreement of March 1921. Failing the abandonment of these activities, and the settlement of a number of minor outstanding claims within ten days, the British Government threatened to annul the trade agreement and withdraw its representative in Moscow. The Soviet Government, frightened by this violent onslaught, agreed to comply with most of the demands, and entered into a mild and inconclusive argument on the propaganda question; and for the time being the storm blew over. 

In Germany, the one major country which had so far accorded de jure recognition to the Soviet Government, the year was also marked by disquieting events. The German economy and the German currency collapsed under the pressure of the Ruhr occupation; and a series of political crises encouraged optimistic observers in Moscow to scent an opportunity to retrieve the failure of March 1921 (see pp. 19, 44 above). In August Brandler and other leaders of the KPD were summoned to Moscow, and plans were laid for a coup to seize power in the autumn. But confidence was sapped by differences about tactics. The whole scheme was bungled in a way which led afterwards to endless recriminations. An isolated communist rising in Hamburg on October 23 was easily crushed. By this time Stresemann was installed as the head of a government pledged to restore the shattered economy; and Seeckt, the head of the Reichswehr, showed every confidence in his ability to maintain order. The paradoxical feature of this episode was that it did not disturb German-Soviet relations. The moral was clear. Seeckt, assured of freedom to deal with German communists, had every incentive to continue and develop military collaboration with Moscow, and Stresemann willingly fell in with this policy. The Soviet Government, in Germany as in Turkey, could not afford to support native communists at the expense of its need for allies and partners in the game of international diplomacy. The same lesson could be drawn from its readiness to cultivate friendly relations with Mussolini's Fascist regime in Italy. 
% 注：去除了段首错位的 "[were", "Id to", 并修复了 "leature" 的 OCR 拼写错误为 "feature"，清理了 "lUU"。

The year 1924 opened under more promising auspices. The advent to power of the first British Labour government brought de jure recognition of the Soviet Government on February 1; and Italian recognition followed a few days later. In May elections in France resulted in the formation of a Left coalition under Herriot. But, owing to the powerful opposition of French holders of pre-revolutionary Russian bonds, recognition of the Soviet Government was delayed till October. During the summer negotiations went on in London for an Anglo-Soviet treaty to replace the trade agreement of 1921. A treaty, accompanied by the promise of a loan, was signed in August in the face of stiff opposition from British financial and commercial interests, and from the Conservative Party. At this point the Liberals withdrew their support from the Labour government, which was defeated in the House of Commons. The treaty was not ratified, and in the ensuing elections the Conservatives won a sweeping victory. Their success was aided by the publication, just before the election, of the ``Zinoviev letter''---a letter of instruction from Comintern to the CPGB to conduct propaganda in the armed forces and elsewhere. The letter was almost certainly a forgery. But its contents seemed plausible; and it was sufficient to inflame public opinion still further against the USSR and its British friends. The new Conservative government, with Austen Chamberlain as its Foreign Secretary, did not formally break off relations, but virtually suspended all dealings with the Soviet Government throughout 1925. Franco-Soviet negotiations for the settlement of debts and claims reached a similar deadlock. 
% 注：修正了 "Febru- ary I" 为 "February 1"，修正了 "tiie face" 为 "the face"，去除了 "I break" 等杂乱符号。

Meanwhile the balance of forces in Europe was changed by the acceptance in August 1924 of the ``Dawes plan'', diplomatically and financially supported from the United States, for an agreed settlement of German reparations obligations with the aid of a massive international loan. This was the starting-point of a process of reconciliation between the victors and vanquished of 1918, which culminated in the famous Locarno treaty, initialled at Locarno in October 1925 and signed with much ceremony in London on December 1. The essence of the treaty was a mutual guarantee of Germany's existing western frontiers---a voluntary acceptance by Germany of this part of the Versailles peace treaty, which did not, however, extend to acceptance of Germany's eastern frontiers. It was ill-received in Moscow, where it was seen as proof of a new westward orientation in German foreign policy and a reversal of Rapallo. It was moreover understood that Germany had been promised admission to the League of Nations with a seat on the League Council; and the Soviet Government expressed particular apprehension that Germany, as a member of the League of Nations, might be obliged to participate in sanctions decreed by the League against the USSR. An attempt was made to meet these fears by a declaration, signed by all the parties to the Locarno treaty, that a member of the League could be required to participate in sanctions only ``to an extent which is compatible with its military situation, and takes its geographical position into account''. On these terms Germany finally entered the League in September 1926. 
% 注：修复了多栏穿插 ("I (obliged to participate... Soviet Government expressed...") 的错误断句，并修正了 "December i." 为 "December 1."。

Assurances to the contrary notwithstanding, the Locarno treaty was rightly assessed in Moscow as an attempt to re-integrate Germany into the western world, to wean her from the Soviet entanglement, and to isolate the USSR as an alien element in the society of nations. The attempt was not wholly successful. Germany, still smarting under the humiliation of 1918, was conscious of an inferior status among the western Powers, and did not wish to become exclusively dependent on them. Association with the USSR was no longer as intimate as in the days when the Rapallo treaty brought the two outcasts together. But it remained for Germany a bargaining counter in relations with the west, and an important factor in the European balance of power. Common mistrust of Poland remained a firm link between the two countries. The secret German-Soviet military arrangements were working well; and the Reichswehr would have strongly resisted anything that tended to disrupt them. Economic relations were profitable to both countries. At the very moment when Stresemann was negotiating with Chamberlain and Briand at Locarno, a German-Soviet trade agreement, carrying with it a substantial credit from a group of German banks, was signed in Moscow. For the USSR Germany was its largest and most reliable trading partner. 
% 注：去除了 "Igrate", "rtrvr" 等跨栏边缘乱码。

This was not, however, the only demonstration of German concern to maintain a foothold in eastern Europe. The Soviet Government, not content with denouncing British efforts to build up an anti-Soviet coalition of states, sought at this time to establish special relations with other states which might be interested to resist this design. But, since it was unwilling to undertake any military commitment, and was concerned primarily to forestall combined action against the USSR, the formula proposed was a mutual undertaking by each party not to participate in hostile action, military or economic, against the other, and to remain neutral in the event of a war arising from aggression against the other. A treaty with Turkey on this basis was signed in December 1925. The same formula, with verbal variants, was embodied in a German-Soviet treaty of April 24, 1926. Some Germans invoked the precedent of Bismarck's ``reinsurance treaty'' with Russia in 1887; and the treaty caused considerable annoyance in the west. Angry interludes from time to time disturbed normal intercourse between Moscow and Berlin. The most serious of these incidents occurred in December 1926, when Soviet shipments of war material to Germany under the secret military agreements came to the knowledge of the German social-democrats, who made a public protest in the Reichstag, to the grave embarrassment both of the German and Soviet Governments, and especially of the German communists and Right-wing nationalists. But fears of Allied reprisals did not materialize; the western Powers were too much involved in maintaining the good relations with Germany established by the Locarno treaty to raise this awkward question. The storm subsided; and during the next few years, while Soviet relations with western Europe were almost a blank, relations with Germany, political, military, economic and cultural, remained far closer and more fruitful than with any other country. 
% 注：去除了段尾处的三个 "I" 侧边栏划线噪点。

The revolutionary element in policy and outlook in the relations of the USSR with the outside world, embodied institutionally in Comintern, still sometimes appeared to conflict with the diplomatic activities conducted by Narkomindel in a way which created momentary embarrassment. But the illusory character of the supposed clash between the claims of revolution and diplomacy, encouraged by the pretence that the Soviet Government had no responsibility for the proceedings of Comintern, was revealed in the argument, constantly repeated, that the USSR was the one solid bulwark of world revolution, whose prospects depended on its strength and security. The interests of international revolution and the national interests of the USSR were on this hypothesis inseparable. A corollary of this view was the dependence of all other communist parties, often referred to as ``sections'' of Comintern, on the Russian party. Any clash between Comintern and the Russian party was unthinkable. When in the spring of 1922 twenty-two members of the Workers' Opposition appealed to Comintern against their treatment by the Russian party, as the statutes entitled them to do, the appeal was rejected out of hand by a commission which included the Bulgarian, Kolarov, and the German, Klara Zetkin. The Russian party alone had led a victorious revolution. It had acquired the right and duty to lead and instruct other parties on the road to revolution. The historical fact that Comintern had developed as an institution built on a Russian model, and centred round the Russian party, lent support to this argument. 
% 注：清除了 "lit tit." 的污渍噪点。

The relation of communist parties to the central organs of Comintern was the key-note of its fifth congress which met in June 1924. The KPD leaders who had bungled the October rising in Germany were condemned as Rightists, and replaced by new leaders from the Left, Ruth Fischer and Maslow. A similar shift occurred in the French and Polish parties, whose leaders, now branded as Rightists, had declared in favour of Trotsky. But, amid much rhetoric at the congress about the virtues of the Left, it was apparent that the main quality demanded of the new Left leaders was disciplined obedience to decisions taken in Moscow. Zinoviev launched the slogan of the ``Bolshevization'' of the parties, defined in a resolution of the congress as ``the transmission to our sections of everything that was and is international, and of general significance, in Russian Bolshevism''. Its adoption seemed a matter of course. It was the automatic product of the delay of revolution in other countries; and it drew fresh reinforcement from the doctrine of socialism in one country, which registered the role of the USSR as the unique exemplar of a successful socialist revolution. Stalin, who had hitherto taken no part in the work of Comintern, modestly attended the fifth congress, but left the limelight to Zinoviev, spoke in some of the commissions, though not in the plenary sessions, and made himself known to foreign delegates. Trotsky, who was present, and drafted a manifesto of the congress on the approaching tenth anniversary of the war of 1914, did not speak. 

For the next three years the isolation of the USSR in the capitalist world was a source of growing anxiety in Moscow. The capitalist economies of Europe, severely shaken by the war of 1914-1918, had by the middle nineteen-twenties regained their equilibrium, and were enjoying a wave of prosperity, stimulated by American investment. Recognition by Comintern that the western countries had achieved a state of ``capitalist stabilization'' was qualified by attaching to it the epithets ``relative'' and ``temporary'', and was matched by insistence on ``Soviet stabilization''. But these considerations inspired a mood of caution. Left leaders of foreign parties who had enjoyed favour at the fifth congress were removed within the next two years, and succeeded by moderates. Annual congresses of Comintern were abandoned, being replaced by ``enlarged'' sessions of IKKI; the sixth congress was not convened till 1928. Visions of the coming revolution were still evoked, but with diminishing conviction. Revolutionary propaganda was conducted, but mainly as a defensive weapon against governments whose hostility was known and feared. The rise of Stalin was greeted with some satisfaction in the west, since it represented the eclipse of revolutionary firebrands like Trotsky and Zinoviev by a moderate and cautious leader, primarily devoted to restoring the fortunes of his own country. 
% 注：修复了 "k capitahst" 为 "the capitalist" 并去除了 "11", "111", "//", "^|||", "i^//f" 等大量的侧边栏识别乱码。

This period was the heyday of the united front, when cooperation of communists with other Left-wing parties and groups was assiduously preached, and of the organization of international ``fronts'', not ostensibly communist, though encouraged and partially financed from Moscow, which recruited Left sympathizers of heterogenous groups or parties to support causes favoured by Comintern. The most famous and successful of these was the League against Imperialism, whose founding congress in Brussels in February 1927 brought together for the first time delegates from China, India and Indonesia, from the Middle East, from many parts of Africa, from Latin America, and from the Negroes in the United States, on a platform of protest against the tyrannical rule of the imperialist Powers over subject peoples. The Moscow celebrations of the tenth anniversary of the revolution in November 1927, attended by a galaxy of distinguished foreign guests, were the occasion for the foundation of an international society of Friends of the Soviet Union. Organizations like the International Workers' Aid and the International Class War Prisoners' Aid, centred in Moscow but with branches in the other principal countries, served the same purpose of maintaining contacts with the non-communist Left and of wooing sympathy for the USSR. 
% 注：修复了由于分栏导致的句子颠倒 ("ill States, on a platform... from Latin America...") 以及去除了 "Uii" 噪点。

Relations with the British labour movement were from the first anomalous. The CPGB had been formed in 1920 by the amalgamation of several splinter groups of the extreme Left; its total membership in the middle nineteen-twenties was about 5000. Its weakness was offset by the unique strength of the British trade unions, which formed the hard core of the workers' movement, and had a dominant influence in the Labour Party. Moreover, the unions had shown on more than one occasion warm and effective sympathy for the Russian revolution and the Soviet regime. To win over the trade unions in capitalist countries was the task of the Red International of Labour Unions (Profintern or RILU) set up in Moscow in 1921. In France and in Czechoslovakia, its efforts were successful in splitting the movement more or less equally between trade unions affiliated to the existing International Federation of Trade Unions (IFTU), commonly called the Amsterdam International, and unions affiliated to Profintern. In Germany no split occurred, and members of the KPD exercised considerable influence in the unions affiliated to Amsterdam. In Great Britain the trade unions, with very few exceptions, remained loyal to Amsterdam. But a majority of British trade unions continued for many years to deplore the split in the international movement, and to call for a reconciliation between the two rival federations. Acute jealousies, as well as deep ideological differences, between Amsterdam and Moscow made this a hopeless ambition. 
% 注：去除了 "jljtbe", "::^l(Red" 等错落的干扰符号。

Profintern was founded at the very moment when Comintern was turning to united front policies. When Lenin at the second congress in 1920 first adumbrated the ideas which took shape a year later under the catchword of the ``united front'' (see p. 16 above), his remarks were directed primarily to British affairs, and to the need for British communists to support the ``MacDonalds and Hendersons'' of the Labour Party, whose peculiar constitution made it possible and normal for members of the CPGB to remain at the same time members of the Labour Party. But in Britain the trade unions provided the most natural ground on which the appeal to sympathizing non-communist workers could be made. The typical British communist was said to carry three membership cards in his pocket: of the CPGB, of his trade union and of the Labour Party. Profintern set up a bureau in London; and, responding to this stimulus, the CPGB promoted two united front organizations---the National Minority Movement (NMM) to act as a ginger group within the trade unions, and the National Unemployed Workers Movement (NUWM) to conduct propaganda and agitation, under communist leadership but with broad workers' participation, on one of the major ills of the period. Though the Labour Party rejected repeated requests from the CPGB for affiliation, its rank-and-file members were not initially inhospitable to individual communists. In the 1922 elections two communists were elected to parliament, one as the official Labour Party candidate, the other with tacit Labour support. 
% 注：去除了 "^(/•from" 和 "^(Knot" 的乱码前缀。

Reaction came more swiftly in the Labour Party than in the trade unions. In 1924 the Labour Party prohibited the selection of communists as official Labour candidates. A decision to ban members of the CPGB from Labour Party congresses was taken, but could not be enforced so long as trade unions included them in their delegations. Sympathy for the USSR in the unions was a hardier growth. Tomsky, the Soviet trade union leader, spoke amid scenes of enthusiasm at the British trade union congresses of 1924 and 1925, and a British delegation attended the Soviet trade union congress in December 1924---just after the Zinoviev letter and the defeat of the Labour government. Early in 1925 a joint Anglo-Russian trade union committee was formed with the aim of fostering cooperation between the unions of the two countries. But the project under-estimated the discrepancies and differences of outlook between Soviet and British trade union leaders, and the unwillingness of the latter to place themselves in opposition to the Amsterdam International. The meetings of the committee were the occasion of increasingly sharp recriminations between Soviet and British delegates. The activities of Profintern, and many tough Soviet criticisms of the British leaders, were resented; and the aggressive tactics of the NMM and the NUWM caused continual irritation. In the General Council of the TUC an anti-Soviet majority confronted a dwindling pro-Soviet minority. 
% 注：去除了 "Ill" 和 "III" 的大写字母边界噪点。

The dividing line was the British general strike of May 1926. In Soviet eyes, a general strike was a political act, a bid for power, an act of class warfare and the beginning of a proletarian revolution. On the British side, it remained, as it had begun, a dispute about wages. The trade union leaders, and a vast majority of the workers, were seeking to extract a fairer share of benefits from the existing system, not to overthrow it. The exhortations to revolution radiated from Moscow alarmed and alienated them; and they refused the financial aid offered by the Soviet unions on the plea that it would prejudice their cause---an insult for which the British leaders were never forgiven by their Soviet counterparts. When after ten days they acknowledged defeat and called off the general strike, leaving the miners, whose wrongs had been its original cause and inspiration, to struggle on helplessly alone, this seemed in Soviet eyes conclusive proof that the British trade union leaders had sold out to the bourgeoisie, and that the only hope lay in raising the rank and file of the workers in revolt against a treacherous trade union bureaucracy. Henceforth Soviet animosity against the British leaders was implacable; and failure to shake the loyalty of the majority of British trade unionists to their leaders embittered and frustrated Soviet relations with the movement for many years. 
% 注：去除了 "__" 以及跨行断裂的 "(the British trade union leaders..." 前置乱码。

The general strike, and the financial aid to the strikers proffered from Moscow, added fuel to the anti-Soviet campaign conducted by prominent Conservative politicians with increasing vehemence ever since the autumn of 1924. During the winter of 1926--1927 the demand in Conservative circles to break off relations with the USSR became irresistible. In May 1927 the premises of Arcos, which contained some offices of the Soviet trade delegation in London, were raided by the police. The seized documents produced no sensational revelations. The purpose of the exercise was, however, clear and was not to be frustrated. On May 24 Baldwin announced the severance of diplomatic relations with the USSR and the annulment of the trade agreement. No other government followed the British example. But the British presence still dominated the European diplomatic scene. The gesture was sufficient to create widespread anxiety. Fears of war, or at the least of an economic and financial blockade, were rife in Moscow. Pilsudski had seized power in Poland in the previous year; and apprehension was felt that Britain might instigate, or support, military action by him against the USSR. The British TUC added to the discomfiture of the Soviet leaders by voting, at its annual congress in September 1927, to dissolve the Anglo-Russian trade union committee, which had long been a target of attack by Trotsky and the opposition in Moscow. No ray of light appeared on the horizon. A grain collections crisis of formidable dimensions followed the harvest. The battle with the opposition in the party reached its peak of bitterness. Even in Asia Soviet fortunes were at their lowest ebb. 
% 注：修复了本段严重的分栏串行错位 ("iand was not to be frustrated. On May 24 Baldwin announced the police... clear (the severance of diplomatic relations") 并去除了诸多侧边标点。

Throughout this period the United States Government inflexibly refused to recognize the Soviet Government or to have any relations with it. This attitude was reiterated by successive Presidents and Secretaries of State, and was challenged only by a handful of radical intellectuals and by a few bankers and businessmen interested in a revival of trade with the USSR. After an official embargo on trade had been removed, a ban on loans to the USSR, a veto on the acceptance of Soviet gold on the alleged ground of contested ownership, and a refusal by the banks to grant credit, constituted an effective prohibition on any large transaction. But a trickle of trade began to flow. In 1924 the Soviet authorities set up a trading corporation in New York under the name Amtorg, the counterpart of Arcos in London. An unofficial Soviet emissary resided in Washington, and occasionally visited the State Department in a private capacity. In 1925 the American financier Harriman secured a concession to work the manganese mines in the Caucasus. Though the project was not a success, and the concession was later cancelled, it was a breaking of the ice. But it was not till after 1927, when industrialization was on the way in the USSR, that American industrialists became seriously interested in the Soviet market. 
% 注：去除了 "i JnI ^ I" 这个极端的 OCR 污渍错误。